% chapters/ch07.tex
\chapter{Conclusiones, aportes e implicaciones de la investigación}

\begin{refsection}


\section{Conclusiones generales de la investigación}

La presente investigación tuvo como propósito diseñar y validar un modelo de diagnóstico educativo inclusivo basado en datos, implementado mediante un sistema tecnológico capaz de apoyar la identificación de necesidades educativas diversas y contribuir a la toma de decisiones pedagógicas fundamentadas.

A partir del análisis teórico, metodológico y empírico desarrollado a lo largo de la investigación, es posible formular varias conclusiones relevantes.

En primer lugar, los resultados evidencian que la integración entre inclusión educativa, Diseño Universal para el Aprendizaje y educación basada en datos constituye un marco conceptual sólido para el desarrollo de herramientas tecnológicas orientadas al diagnóstico educativo. Esta articulación permite superar enfoques tradicionales centrados exclusivamente en el rendimiento académico y promueve una comprensión más amplia y contextualizada del aprendizaje.

En segundo lugar, el modelo de diagnóstico educativo inclusivo propuesto demuestra que la analítica de datos educativos puede contribuir a identificar patrones de aprendizaje relevantes, facilitando la comprensión de las diferencias individuales entre los estudiantes. La utilización de múltiples fuentes de información permite construir perfiles de aprendizaje más completos y contextualizados.

En tercer lugar, la implementación del sistema tecnológico en un contexto educativo real evidenció que las herramientas digitales basadas en datos pueden desempeñar un papel significativo en el apoyo a la toma de decisiones pedagógicas. Los docentes participantes señalaron que la información generada por el sistema les permitió comprender mejor las características de aprendizaje de sus estudiantes y ajustar sus estrategias didácticas en consecuencia.

Finalmente, la investigación confirma que el uso de tecnologías educativas basadas en datos debe desarrollarse bajo principios éticos claros, incluyendo la protección de la privacidad de los estudiantes, la transparencia en el uso de los datos y la prevención de sesgos algorítmicos. Estos principios resultan fundamentales para garantizar que el uso de la tecnología contribuya efectivamente a la inclusión educativa.

\section{Aportes científicos de la investigación}

La investigación aporta al campo de la educación inclusiva y la tecnología educativa en distintos niveles.

\subsection{Aporte teórico}

Desde el punto de vista teórico, la investigación contribuye a la articulación entre tres enfoques relevantes en la educación contemporánea:

\begin{itemize}
\item la inclusión educativa,
\item el Diseño Universal para el Aprendizaje,
\item la educación basada en datos.
\end{itemize}

Esta articulación permite ampliar el marco conceptual del diagnóstico educativo, integrando perspectivas pedagógicas, tecnológicas y metodológicas que tradicionalmente han sido abordadas de manera separada.

Asimismo, el estudio aporta a la reflexión sobre el papel de los sistemas tecnológicos inteligentes como herramientas de mediación pedagógica orientadas a apoyar los procesos educativos.

\subsection{Aporte metodológico}

En el plano metodológico, la investigación demuestra la pertinencia del enfoque de Investigación Basada en el Diseño (\DBR) para el desarrollo y evaluación de innovaciones tecnológicas aplicadas a la educación.

La aplicación del enfoque \DBR\ permitió:

\begin{itemize}
\item diseñar un modelo conceptual fundamentado en el marco teórico;
\item desarrollar un sistema tecnológico prototipo;
\item evaluar su funcionamiento en un contexto educativo real;
\item generar conocimiento teórico y práctico a partir de la experiencia de implementación.
\end{itemize}

Este proceso metodológico contribuye al fortalecimiento de las metodologías de investigación orientadas al estudio de tecnologías educativas.

\subsection{Aporte tecnológico}

La investigación también aporta al desarrollo de herramientas tecnológicas orientadas al diagnóstico educativo inclusivo.

El sistema desarrollado permite:

\begin{itemize}
\item recopilar datos educativos de manera sistemática;
\item analizar patrones de aprendizaje mediante herramientas de analítica educativa;
\item generar diagnósticos pedagógicos orientados a la inclusión educativa.
\end{itemize}

Aunque el sistema se presenta como un prototipo académico, su diseño modular permite futuras ampliaciones y adaptaciones en distintos contextos educativos.

\subsection{Aporte educativo y social}

Desde una perspectiva educativa y social, la investigación contribuye a promover prácticas pedagógicas más equitativas mediante el uso responsable de tecnologías basadas en datos.

El modelo propuesto permite reconocer la diversidad de los estudiantes como un elemento central del proceso educativo y ofrece herramientas para apoyar la toma de decisiones pedagógicas orientadas a la inclusión.

\section{Implicaciones para la práctica educativa}

Los resultados de la investigación sugieren diversas implicaciones para la práctica educativa.

En primer lugar, el uso de sistemas tecnológicos basados en datos puede contribuir a fortalecer los procesos de diagnóstico educativo, proporcionando a los docentes información más detallada sobre las características de aprendizaje de sus estudiantes.

En segundo lugar, la integración entre tecnología educativa y principios pedagógicos como el Diseño Universal para el Aprendizaje permite desarrollar estrategias de enseñanza más flexibles y adaptadas a la diversidad del alumnado.

En tercer lugar, el uso de analítica de datos educativos puede facilitar la toma de decisiones pedagógicas fundamentadas, siempre que estas herramientas se utilicen como apoyo al juicio profesional del docente y no como sustituto de su experiencia pedagógica.

\section{Limitaciones de la investigación}

Como toda investigación aplicada, el presente estudio presenta ciertas limitaciones que deben ser consideradas en la interpretación de los resultados.

Entre las principales limitaciones se encuentran:

\begin{itemize}
\item el tamaño limitado de la muestra utilizada en la implementación del sistema;
\item el carácter contextualizado del estudio;
\item la naturaleza prototípica del sistema tecnológico desarrollado.
\end{itemize}

Estas limitaciones no invalidan los resultados obtenidos, pero sugieren la necesidad de continuar investigando el modelo propuesto en otros contextos educativos y con muestras más amplias.

\section{Líneas futuras de investigación}

A partir de los resultados obtenidos en la investigación, se identifican diversas líneas de investigación futura que podrían ampliar y profundizar el modelo propuesto.

Entre ellas se destacan:

\begin{itemize}
\item la aplicación del modelo de diagnóstico educativo inclusivo en distintos contextos educativos;
\item el desarrollo de nuevas funcionalidades basadas en inteligencia artificial para mejorar los procesos de análisis de datos educativos;
\item la integración del sistema con plataformas educativas existentes;
\item el análisis longitudinal del impacto del modelo en los procesos de aprendizaje.
\end{itemize}

Estas líneas de investigación pueden contribuir a fortalecer el desarrollo de tecnologías educativas orientadas a la inclusión y a la mejora de los procesos de enseñanza y aprendizaje.

\section{Reflexión final}

La educación contemporánea enfrenta el desafío de atender la diversidad de los estudiantes en contextos educativos cada vez más complejos. En este escenario, el uso responsable de tecnologías basadas en datos puede ofrecer nuevas oportunidades para comprender los procesos de aprendizaje y apoyar la toma de decisiones pedagógicas.

Sin embargo, la tecnología por sí sola no garantiza la inclusión educativa. Su efectividad depende de la manera en que se integre con principios pedagógicos sólidos, prácticas docentes reflexivas y marcos éticos claros.

En este sentido, la presente investigación propone un modelo que busca contribuir a este desafío, integrando inclusión educativa, tecnología y analítica de datos en un enfoque orientado a fortalecer las prácticas educativas inclusivas.